\begin{Definition}{Probability space $(\Omega, \mathcal{F}, P)$}
    -Donde $\mathcal{F}$ es una sigma algebra de $\Omega$.
\end{Definition}

\begin{Definition}{Función de probabilidad}
    -$Pr: \mathcal{F} \to \reals$ que satisface que:
    $$ \forall E \in \mathcal{F} : P(E) \in [0, 1]$$
    $$ Pr(\Omega) = 1 $$
    $$ (E_n)_{n \leq \naturals} : E_n \cap E_{n+1} = \emptyset \then P(\cup_{i \leq \naturals} E_i) = \sum_{i \leq \naturals} P(E_i) $$
\end{Definition}

\begin{Lemma}{Lema de la union}
    - $E_1, E_2$ eventos arbitrarios:
    $$ \then P(E_1 \cup E_2) = P(E_1) + P(E_2) - P(E_1 \cap E_2) $$
\end{Lemma}

\begin{Lemma}{Lema de la union extendida}
    - $E_1, E_2, \dots, E_n$ eventos arbitrarios:
    $$ P(\cup_{i =1}^N E_i) = \sum_{i=1}^n P(E_i) - \sum_{(i, j) tuplas combinaciones de \{1, \dots, n\}}  + \sum_{(i, j, k) tuplas combinaciones de \{1, \dots, n\}}
    \dots $$
\end{Lemma}

Con este lema se puede demostrar la monoticidad de la probabilidad para $I, J: I \subseteq J$
Con esto a su vez se puede demostrar que:

\begin{Lemma}{Lema de la subaditividad}
    - $$Pr(\cup_{i \geq 1} E_i) \leq \sum_{i \geq 1} P(E_i)$$
\end{Lemma}

\begin{Lemma}{Monotonía de la probabilidad}
    - $$ A \subset B \implies P(A) \leq P(B)$$ \\
    $$ P(\cup_{m = n}^n A_m) \implies P(\cup_{m = n}^n A_m) \leq P(A_k) \quad \forall k \geq n \implies P(\cup_{m = n}^n A_m) \leq \liminf P(A_k) $$
\end{Lemma}

\begin{Lemma}{Eventos independientes}
    -Decimos que $E, F$ son dos eventos independientes ssi:
    $$P(E \cap F) = P(F) \cdot P(E) $$
    Decimos que una coleccion $E_{n: 1\leq n \leq N}$ es mutuamente 
    independiente ssi para cualquier forma de extraer elementos de $\{1, \dots, n\}$ entonces:
    $$ P(\cap_{i \in I} E_i) = \prod_{i \in I} P(E_i)$$
\end{Lemma}

\begin{Definition}{lim inf y lim sup}
    - Definimos lim inf sobre $(A_n)_{n \in \naturals}$ como:
    $${liminf}_{n \to \infty} (A_n) = \lim_{n \to \infty} \inf_{k \geq n} A_n = \cup_{n}^\infty \cap_{k \geq n}^\infty A_n $$
\end{Definition} 
Lim sup se define alternativamente.

\begin{Lemma}{Continuidad de la probabilidad}
    - Sea $(B_n)_{n \in \naturals}$ : $B_n \subseteq B_{n+1}$. Definimos $B = \lim_{n \to \infty} B_n = \cup_{n=1}^\infty B_n$.
    $$\then P(\lim_{n \to \infty} B_n) = P(B) = \lim_{n \to \infty} P(B_n)$$
\end{Lemma}

\begin{Lemma}{Lemma de Fatou debil}
    - $P(\liminf_{n \to \infty} A_n) \leq \liminf_{n \to \infty} P(A_n)$
\end{Lemma}

\begin{Definition}{Igualdad puntual fuerte}
    - $X, Y: \Omega \to \reals $ son iguales si:
    $$ X(w) = Y(w) \quad \forall w \in \Omega $$
\end{Definition}

\begin{Definition}{Igualdad casi seguramente}
    - $X, Y \; c.s \; si \; \mathbb{P}(X \neq Y) = 0 $ \\
    Que es equivalente a decir:
    $ \mathbb{P} ( \{ w: X(w) = Y(w)\}) = 1$
\end{Definition}


\begin{Definition}{Igualdad en distribucion}
    - $X =_d Y \iff \mathbb{P}(X \leq t) = \mathbb{P}(Y \leq t) \forall t $
\end{Definition}

\begin{Definition}{Densidad de X}
    - (f): $f$ es densidad de $X$ v.a entonces $f$ es tal que:
    $$ \mathbb{P}(X \in A) = \int_A f(x) dx $$
\end{Definition}

\begin{Proposition}{Sobre las densidades}
    - Sea $f$ densidad de $X$. Luego:
    $$ \bar{f} = \begin{cases}
        f \; en \; X - \{ w_0\} \\
        f + \epsilon \; en w_0
    \end{cases}$$
    entonces $\bar{f}$ sigue siendo densidad.
\end{Proposition}

\begin{Proposition}{Sobre las densidades}
    -El conjunto de las densidades de $X$ es no numerable.
\end{Proposition}

\begin{Theorem}{Ley debil de los grandes numeros}
    - Sea $\{ X_i \}_{i=1}^N$ v.a iid tal que $\mathbb{E}(X_i) = \mu \quad \forall i \in \{ 1 \dots N \}$
    entonces:
    $$ \frac{S_N}{N} \rightarrow \mu $$
    Casi seguramente. Es decir:
    $$ \forall \epsilon > 0 \exists N_0 \in \naturals : \forall n \geq N_0 : \mathbb{P}(| \frac{S_N}{N} - \mu | \geq \epsilon) \rightarrow 0$$
\end{Theorem}

\begin{Definition}{Propiedad tipo Hoeffding}
    - Sea $(X_i)_{i=1}^N$ una coleccion de v.a iid. Sea $\vec{a} \in \reals^N$. Decimos que hay Hoeffding si:
    $$ \exists c > 0 \forall \vec{a} \in \reals^N: \mathbb{P}(\sum_{i=1}^{N} a_i X_i \geq t )  \leq exp( \frac{-ct^2}{\sum_{i=1}^{N} a_i^2})$$
\end{Definition}

\begin{Proposition}{Hoeffding y la media}
    - $ (X_i)_{i=1}^N $ v.a iid tal que $ \mathbb{E}(X_i) = \mu $ y que satisface Hoeffding
    $\then \mathbb{E}(X_i) = 0$
\end{Proposition}